\section*{Erd\H{o}s Problem 408: an additive formula for totient height}
\subsection*{Statement and outcome}
Set $f(1)=0$ and let $f(n)$ be the number of iterations of Euler's totient needed to reach $1$. Does $f(n)/\log n$ have a limiting distribution or a constant normal order? What happens to the largest prime factor after about $\log\log n$ iterations?

We reconstruct the exact additive structure of the height, prove two-sided logarithmic bounds, and exhibit different limiting ratios on two sparse subsequences. These do not settle the distribution or normal-order questions. Source: \url{https://www.erdosproblems.com/408}. The additive structure is classical Shapiro theory; a primary modern treatment is Bal--Bhatnagar, Theorem 1, \url{https://arxiv.org/abs/1903.09619}. No novelty is claimed.

\subsection*{Constructing the additive function without circularity}
Define a weight $w(p)$ on primes recursively by
\[
 w(2)=1,\qquad w(p)=\sum_{q^e\parallel p-1}e\,w(q)\quad(p>2).
\]
Every prime on the right is smaller than $p$, so this definition is well-founded. Extend completely additively to $h(1)=0$ and
$h(n)=\sum_{p^e\parallel n}e w(p)$. The totient product formula implies
\[
 h(\varphi(n))=
 \begin{cases}
 h(n)-1,&n\text{ even},\\
 h(n),&n>1\text{ odd}.
 \end{cases} \tag{1}
\]
For each odd prime factor $p$, its contribution is $(e-1)w(p)+h(p-1)=e w(p)$; the prime $2$, if present, contributes one less than its original exponent. This proves (1) directly.

For $n>2$, $\varphi(n)$ is even, since reduced residue classes pair as $u,n-u$. Induction on $n$, using $\varphi(n)<n$ for $n>1$ and checking $n=2$ separately, now gives
\[
 f(n)=\begin{cases}
 h(n),&n\text{ even},\\
 h(n)+1,&n>1\text{ odd}.
 \end{cases} \tag{2}
\]
Thus the prime-factor formula for the height is fully proved, rather than assuming additivity of $\varphi$ or of its iterates. In particular for $a,b>1$,
\[
 f(ab)=\begin{cases}
 f(a)+f(b),&a,b\text{ both even},\\
 f(a)+f(b)-1,&\text{otherwise}.
 \end{cases} \tag{3}
\]
The exclusion of $1$ in (3) is necessary with the convention $f(1)=0$.

\subsection*{Bounds and exactly computed extremal families}
We show inductively on primes that $p\le3^{w(p)}$. It holds for $2$. If it holds for smaller primes, then the even integer $p-1$ satisfies
\[
 p-1\le\frac23\,3^{h(p-1)}=2\cdot3^{w(p)-1};
\]
one factor $2$ gives the extra $2/3$, and the other prime factors use the induction hypothesis. Since $w(p)\ge1$, adding one proves $p\le3^{w(p)}$. Multiplication then gives $n\le3^{h(n)}$ for every $n$, and $n\le(2/3)3^{h(n)}$ for even $n$. Equation (2) yields
\[
 f(n)>\log_3 n\quad(n>1). \tag{4}
\]
For odd $n$ use $f=h+1$, and for even $n$ use the strict factor $2/3$.

For an even integer $m$, $\varphi(m)\le m/2$. After the first step from any $n>2$, all subsequent iterates greater than one are even. Hence before reaching one, $\varphi_j(n)\le(n-1)/2^{j-1}$ for $j\ge1$. Taking $j=\lceil\log_2 n\rceil$ forces the iterate below $2$, unless it already reached $1$. Thus
\[
 f(n)\le\lceil\log_2 n\rceil. \tag{5}
\]
The rounding and the initial parity step matter: $\varphi(m)\le m/2$ is false for odd primes.

As $w(3)=h(2)=1$, the exact formulas include
\[
 f(2^r)=r,\qquad f(3^r)=r+1,\qquad f(2\cdot3^r)=r+1.
\]
Therefore the ratios tend respectively to $1/\log2$ and $1/\log3$ on the first and third families. There is no ordinary limit over all integers. These families have density zero, so they do not contradict a constant normal order or a limiting distribution.

\subsection*{A modest prime-factor observation and the remaining gap}
Let $P(m)$ be the largest prime factor, with $P(1)=1$. Every prime dividing $\varphi(m)$ either is a prime already dividing $m$, or divides $p-1$ for some prime $p\mid m$. Consequently $P(\varphi(m))\le P(m)$, and the largest prime factors never increase under iteration.
For $1\le j\le f(n)$, the halving estimate also gives
$P(\varphi_j(n))\le(n-1)/2^{j-1}$. At $j=\lfloor\log\log n\rfloor$ this is only $O(n/(\log n)^{\log2})$, far larger than $n^{o(1)}$. Monotonicity alone does not provide the stronger typical bound requested.

The older local draft's valid elementary upper bound is retained, now joined by the full additive reconstruction and its consequences. Understanding the weighted prime-factor sum $h(n)$ for a density-one set of integers, and controlling large prime factors along typical totient chains, remain unresolved in this attempt. No conjectural distribution theorem is used as an unconditional result.
\textbf{Completion rate estimate: 45\%.}
